\section{Cycle--Overlap Visibility Lemma (Empirical)}

\subsection*{Statement}

Fix $\Delta$ and $R$.  
There exists a threshold $T$ such that for any graph $G$ with
\[
\maxdeg(G)\le \Delta
\]
and
\[
COR_R(G)\ge T,
\]
there exist vertices $u,v$ satisfying
\[
tp_3(u)\ne tp_3(v).
\]

Equivalently,
\[
\chi_{WL^2}(u)\ne \chi_{WL^2}(v).
\]

\subsection*{Experimental Evidence}

Graph: rr\_n12000\_d4\_seed9

Radius:
\[
R=6
\]

Measured:

\[
COR_R(G)=649
\]

Overlap components:
\[
253
\]

Largest component:
\[
82
\]

\paragraph{WL Tests}

\[
WL^1:\ \text{homogeneous}
\]

\[
WL^2:\ \text{distinct colors} = 7
\]

Hence

\[
\exists u,v:\chi_{WL^2}(u)\ne\chi_{WL^2}(v).
\]

\subsection*{Interpretation}

Cycle–overlap rank becomes visible at

\[
FO^3
\]

while remaining invisible to

\[
FO^1,\ FO^2.
\]

\subsection*{Role in the Oblivion Atom}

Provides empirical support for the bridge

\[
COR_R(G)\ \Rightarrow\ FO^k\text{-type diversity}
\]

with

\[
k=3.
\]

